\chapter{}

I compaesani di Domenico, quando andavano a Siena, mangiavano sempre alla sua trattoria; portandogli i saluti e le notizie dei parenti, e magari una fazzolettata di frutta.

Uno di costoro, volendo che il suo figliolo Antonio imparasse a fare il muratore, come a Civitella non avrebbe potuto, gli chiese che lo affidasse e lo raccomandasse a qualche bravo capomastro. Domenico, i giorni di festa, lo invitava a stare con Pietro; e così ambedue i giovanetti, ch’erano quasi della stessa età, dovettero doventare amici, sebbene non andassero d’accordo; ed Agostino, che aveva antipatia per Antonio, fu sostituito.

E siccome, per passeggiata, soli, arrivavano quasi sempre, come voleva il trattore, a Poggio a’ Meli, dopo qualche mese Antonio si vantò di aver parlato di nascosto con Ghìsola. Ed era vero; ma Pietro, da prima, suppose che mentisse, con una delusione violenta, con un dispiacere che pigliava tutto il suo amor proprio. Un amico non doveva mentire. Che aveva detto a Ghìsola? E perchè le aveva parlato senza avvertirlo?

Quale umiliazione provava quando gli altri non rispettavano i suoi sentimenti e obbligavano la sua anima a disfarsi!

Gli altri facevano di lui quello che volevano, e a lui si stringeva la gola dall’emozione. Arrossiva, si sgomentava; sentivasi perso. E nessuna cosa era adatta per lui: le strade troppo faticose, il sole troppo caldo, gli abiti tagliati male, le mani troppo grosse; affannandosi a non riflettere a ciò, di convincersi del contrario; stordendosi; mentre gli orecchi gli rombavano, e credeva di dover cadere da un momento all’altro.

Gli sembrava che la sua faccia non fosse capace a nascondere la lealtà troppo aperta e ostinata; provandone una violenza che gli dava il malessere. Si sentiva debole sotto il suo spirito affannato, che egli stesso voleva cambiare.

Una domenica, tra le altre, tornò con Antonio a Poggio a’ Meli; perchè aveva scommesso di farlo passare da bugiardo dinanzi a Ghìsola. Ma si vergognava di dirgli quel che soffriva dentro di sè; e sentivasi così da meno del suo amico che gli pareva di statura anche più alta del solito.

Già, camminando, s’erano bisticciati, picchiandosi su la schiena; ed egli aveva piuttosto voglia di smettere e di piangere, disperato che l’altro, invece, ci si divertisse.

Antonio, avvedendosi facilmente del turbamento di Pietro, gli gridò:

— Vedrai se non è vero!

Pietro non rispose più; e l’amico soggiunse:

— Le parlai anche l’altro giorno. Ha promesso di voler bene a me e non a te.

E, per troncar corto, gli dette un pugno; ma Pietro se lo riparò con una mano.

Antonio, sempre più sicuro, seguitava a ripetere:

— Tu non ti avvicinerai a lei.

— Nè meno tu.

— Io farò quello che voglio.

E fingendosi risentito, si riaccostò con la saliva bianca che gli usciva di bocca. Anche quando non parlava gli si vedevano tutti i denti di sopra, sani ma storti: sembrava che li avesse piantati nel labbro. E aveva il naso piegato da una parte.

Pietro, cercando di persuaderlo con la bontà, gli disse:

— Ed io mi adirerò con te.

— E che m'importa? Fai quello che vuoi. Io sono amico di tuo padre, e verrò quando mi pare. Anzi tuo padre, qui al podere, mi ci porta più volentieri che te.

Pietro si sentì combattuto senza riparo: era proprio vero quel che aveva detto!

E seguitarono a camminare accanto. Ma, dopo un poco, Antonio lo fermò per guardarlo in faccia; trattenendolo per un braccio. Poi fece una sghignazzata:

— Stai zitto?

Poi sputò sull'erba, asciugandosi la bocca con il dorso della mano.

Pietro disse:

— Io torno indietro.

— Io no: voglio parlarci. Vattene.

— Torna indietro anche tu.

Voleva evitare che Antonio la vedesse. Ma quegli proseguiva; e, allora, Pietro dovette fare altrettanto.

Quando giunsero davanti all'aia, Ghìsola usciva di casa proprio in quel mentre; e s'avviava nel campo a chiamare il nonno, passando accanto alla bella pianta di ciliegio da capo a un filare di viti.

Antonio, per fare il più bravo, le mosse incontro in fretta. Ma Ghìsola rise di più a Pietro; e dette a capire che si fermava lì per lui.

Allora Antonio si mosse per cogliersi una piccia di ciliegie, lasciandoli discosti; e Pietro le domandò:

— È vero che vuoi bene soltanto a me? Dimmelo. Se non fosse vero....

Gli rispose con dolcezza:

— Soltanto a lei.... Però, Antonio non vorrebbe.

Allora non si sentì sicuro, e guardò il dorso dell'amico.

Ghìsola, accortasene, aggiunse:

— Non ci crede?

E scosse la testa. Ella parlava, questa volta, con una tranquillità così profonda, ch'egli fu subito rassicurato.

— Ma non se ne faccia accorgere da lui. Perchè ce lo porta?

Gli sembrò che lo rimproverasse di non stare a solo con lei e credette che ne soffrisse.

Ma la sua bellezza lo distrasse e gli fece dimenticare quel che Antonio aveva detto.

Antonio, intanto, si riavvicinò; certo dopo aver progettato qualche cosa, sputando lontano i noccioli delle ciliegie mangiate tutte insieme; aiutandosi con un dito per cacciarseli di bocca. Pietro, mentre un brivido lo scuoteva, gliene strappò una piccia infilata alle dita. Antonio esclamò:

— Perchè me le levi? Dàlle a Ghìsola, piuttosto.

Pietro non seppe che rispondere; perchè avrebbe voluto che quella cosa non gli fosse stata suggerita; e restò con le ciliegie in mano. Ma Ghìsola lo cavò d'impaccio:

— Io le prendo da me.

Quanto gli parve buona e intelligente!

Ma Antonio non si perse d'animo:

— Se non ci arrivi, ti abbasso il ramo io.

Allora, Pietro notò come a lui non sfuggiva mai nulla per ingraziarsela; ma Ghìsola, aspettandosi anche questo, sorrise e disse:

— Non importa.

Ma con una insolenza, che Pietro sussultò sorpreso. E pensò: «Perchè non è venuto a me di dirglielo prima? Ora non c'è più tempo! E quanto piacere ella avrebbe avuto se glielo avessi detto io!».

Si guardarono tutti e tre in silenzio, stando in cerchio; ma si sentirono per un istante amici e senza ostilità. E sentirono anche il bisogno di dirsi più di quello che s'erano detto fino ad allora.

Ghìsola sembrava più lieta, si mandava in dietro i capelli; toccava il laccio del grembiule, come per invitare a farselo sciogliere. Ma Pietro credeva che se ne volesse andare, perchè non riesciva a dirle niente.

Il ciliegio aveva il pedano nero e rossiccio, aperto da profonde screpolature come spacchi, ripieni di resina dura e lucente; una fila di formiche, saliva, ed un'altra, accanto, scendeva, brulicanti; pareva di sentirsele camminare addosso. Vicino, su l'erba acciaccata, c'era rimasta una pozzanghera di solfato di rame incalcinato. Sopra un fragolaio pendeva un fico, senza nè meno una foglia, tutto liscio, con i rami quasi arruffati insieme; e la sua buccia era di un bianco roseo. Qualche rospo s'udiva dai fondi dei borri, tra i salci potati e rossi. Pareva che non ci fosse nessun'ombra; ma le nebbioline, che restavano basse come le piante, salivano dalle terre vangate.

Antonio, vedendo Pietro assorto, lo urtò. Quegli per non cadere fece un passo innanzi, presso Ghìsola; ma non fiatò perchè Antonio non volesse picchiarlo proprio lì: gli parve che ella odorasse molto, di un odore strano; che lo eccitò. Gli parve anche che facesse l'atto di aprirgli le braccia; e ne fu tutto sconvolto: «Se l'avesse aperte da vero?».

Ma Antonio disse a Ghìsola:

— È possibile che tu pensi a lui? Non vedi com'è brutto?

La contadina, specie per rispetto, rispose che non era vero; ma in modo che Antonio non se la prendesse troppo. Poi seguitò a difenderlo:

— Che gliene importa?

Allora Pietro fu quasi sicuro di non essere solo; ma non ebbe la forza d'alzare gli occhi, benché Antonio non sapesse più quel che dire. Poi Pietro la guardò; ed ella gli sorrise con uno di quei sorrisi involontariamente dolcissimi.

Perciò Antonio, non trovando da proporre di meglio, perchè quei due non stessero troppo insieme, disse con tutta la sua cattiveria:

— Io me ne torno a Siena.

Ghìsola suggerì sottovoce a Pietro, sapendo che Antonio avrebbe udito lo stesso:

— Lo lasci andare.

E allora Antonio, senza aspettarlo, s'avviò; ma, volgendosi con collera, chiese:

— E tu non vieni?

Ghìsola non parlava più: e il suo silenzio non lasciava trapelar nulla. Si capiva bene però che voleva mettere alla prova Pietro, che le disse con la voce strozzata:

— Bisogna che vada. Mio padre....

Tutta la faccia di lei s'indurì; ed ella si mise a guardare Antonio già discosto parecchi passi.

Pietro si raccomandò:

— Non dirgli niente!

Ella abbassò la testa, rispondendo:

— Allora vada via!

Ma Pietro credette d'essere amato. E raggiunse Antonio, prendendolo a braccetto. Cominciarono, allora, a ridacchiare.

Poi, Antonio disse sinceramente, e anche perchè Pietro non pensasse più a Ghìsola:

— Perchè siamo venuti a Poggio a' Meli? Non ci siamo divertiti.

Una cicala cantò da un olivo. La saggina ondeggiava prima lenta e poi in fretta; talvolta qualche stelo pareva scosso da un brivido, aprendo a tratti i suoi fiori chiari.

Antonio cavò di tasca un coltellino con il manico d'osso a coda di pesce, spingendolo sotto la buccia secca di una canna, che aveva raccolta; tagliando anche i cerchietti dei nodi, a colpi che assomigliavano al suo riso.

Pietro non si volse indietro a vedere dove fosse Ghìsola perchè non facesse altrettanto Antonio, giacché ora fingeva d'essere attento al suo lavoro di pulitura. Antonio infatti lo spiava; ma era sicuro che non ce ne fosse bisogno.

Giunti alla Porta Camollia, si spolverarono con il fazzoletto le scarpe, si asciugarono il sudore e si ravversarono il cappello aiutandosi a rifarci la piega nel mezzo.

Prima d'entrare nella trattoria, si promisero di non parlare più nessuno dei due a Ghìsola.
